% ============================================================
% Section 119: InP 半导体低温吸收光谱与激子效应
% ============================================================
\section{半导体物理：InP 低温吸收光谱与激子效应}
\label{sec:inp-exciton}

\begin{modelbox}[题目：InP 直接带隙半导体的吸收光谱分析]
直接带隙纯净半导体 InP 在低温下（$T\approx4\,\text{K}$）的吸收光谱中，连续吸收阈值 $E_A=1.400\,\text{eV}$，激子吸收峰 $E_C=1.386\,\text{eV}$。InP 的静态介电常数 $\varepsilon=9.6$，空穴远比电子轻（$m_h^*\ll m_e^*$）。

(1) $A$ 处连续吸收谱和 $C$ 处吸收峰的物理来源是什么？

(2) 计算 InP 中的直接带隙宽度；

(3) 计算空穴的有效质量 $m_h^*$；

(4) 在室温下 $C$ 处的峰观察不清楚，解释这种效应。
\end{modelbox}

\begin{tipbox}[考点快速分析]
\begin{itemize}
    \item \textbf{本征吸收边}：光子能量 $\ge E_g$ 时产生自由电子-空穴对，直接带隙跃迁无需声子
    \item \textbf{激子吸收}：电子-空穴库仑束缚态，峰位在 $E_g-E_{\text{ex}}$，类氢模型
    \item \textbf{激子结合能}：$E_{\text{ex}}=\mu^*e^4/[2(4\pi\varepsilon_0\varepsilon)^2\hbar^2]$，$m_h^*\ll m_e^*$ 时 $\mu^*\approx m_h^*$
    \item \textbf{室温展宽}：声子散射 + 热离解（$k_BT\gtrsim E_{\text{ex}}$）
\end{itemize}
\end{tipbox}

\begin{derivationbox}[标准解题过程]
\textbf{(1) 吸收谱的物理来源}

\textbf{$A$ 处的连续吸收谱（本征吸收边）}：当入射光子能量 $h\nu\ge E_g$ 时，价带电子被激发到导带，产生自由电子-空穴对。InP 为直接带隙半导体，跃迁无需声子参与，吸收边陡峭。阈值 $E_A$ 对应禁带宽度 $E_g$。

\textbf{$C$ 处的吸收峰（激子吸收）}：当 $h\nu<E_g$ 时，光子虽不足以产生自由电子-空穴对，但可激发电子-空穴束缚态——\textbf{激子}。激子能级位于禁带中靠近导带底处。最低能量的激子态（基态激子）对应尖锐的吸收峰 $C$。峰 $C$ 与吸收边 $A$ 之间的能量差即为激子结合能 $E_{\text{ex}}$。

\textbf{(2) 直接带隙宽度}

本征吸收边阈值直接给出禁带宽度：
\begin{equation}
\boxed{E_g=E_A=1.400\,\text{eV}.}
\end{equation}

\textbf{(3) 空穴有效质量的计算}

激子可类比为类氢原子，其结合能（SI 单位制）为
\begin{equation}
E_{\text{ex}}=\frac{\mu^*e^4}{2(4\pi\varepsilon_0\varepsilon)^2\hbar^2},
\end{equation}
其中折合有效质量
\begin{equation}
\frac{1}{\mu^*}=\frac{1}{m_e^*}+\frac{1}{m_h^*}.
\end{equation}

由于 $m_h^*\ll m_e^*$，有 $\mu^*\approx m_h^*$。

激子结合能为吸收边与激子峰的能量差：
\begin{equation}
E_{\text{ex}}=E_A-E_C=1.400-1.386=0.014\,\text{eV}=2.24\times10^{-21}\,\text{J}.
\end{equation}

由结合能公式解出 $m_h^*$：
\begin{equation}
m_h^*=\frac{2(4\pi\varepsilon_0\varepsilon)^2\hbar^2E_{\text{ex}}}{e^4}.
\end{equation}

代入 $\varepsilon=9.6$，$\varepsilon_0=8.854\times10^{-12}\,\text{F/m}$，$\hbar=1.055\times10^{-34}\,\text{J}\cdot\text{s}$，$e=1.602\times10^{-19}\,\text{C}$：
\begin{equation}
m_h^*\approx8.67\times10^{-32}\,\text{kg}\approx0.095\,m_0.
\end{equation}

\begin{equation}
\boxed{m_h^*\approx0.095\,m_0.}
\end{equation}

\textbf{(4) 室温下 $C$ 峰模糊的原因}

室温下晶格热振动（声子）显著增强。激子与声子发生强烈散射，导致能级展宽。此外，热激发能量 $k_BT\approx26\,\text{meV}$（室温）已大于激子结合能（$14\,\text{meV}$），激子极易被热离解为自由电子和空穴。因此激子寿命极短，吸收峰显著展宽并弱化，仅表现为本征吸收边的拖尾或肩部，无法观测到尖锐的 $C$ 峰。
\end{derivationbox}

\begin{formulabox}[答案汇总]
\begin{center}
\begin{tabular}{ll}
\toprule
\textbf{问题} & \textbf{答案} \\
\midrule
(1) $A$ 处连续吸收 & 本征吸收（带间跃迁），阈值 $E_A=E_g$ \\
(1) $C$ 处吸收峰 & 基态激子吸收，电子-空穴束缚态 \\
(2) 直接带隙 $E_g$ & $1.400\,\text{eV}$ \\
(3) 空穴有效质量 $m_h^*$ & $\approx0.095\,m_0$ \\
(4) 室温 $C$ 峰模糊 & 声子散射和热离解导致激子能级展宽、寿命缩短 \\
\bottomrule
\end{tabular}
\end{center}
\end{formulabox}

\begin{insightbox}[物理图像]
\begin{itemize}
    \item 低温吸收谱是研究半导体能带结构和激子效应的有力工具。激子峰的出现是直接带隙半导体低温吸收谱的典型特征，反映电子-空穴库仑相互作用的束缚效应。
    \item 通过激子结合能和类氢模型可估算折合有效质量。本题因 $m_h^*\ll m_e^*$，折合质量主要由轻的空穴决定，从而直接获得 $m_h^*$。
    \item 激子结合能通常较小（$\sim\text{meV}$ 量级），因此室温下易被热离解（$k_BT\gtrsim E_{\text{ex}}$）。激子效应在低温下才显著，是光电器件常在低温下工作的原因之一。
\end{itemize}
\end{insightbox}
